% GENERATED FILE. Edit canonical lessons, then run npm run generate:books.
% canonical-source-hash: fnv1a64:08298a1fb309a467
% canonical-lessons: HI-C247-resam, HI-C247-camra, HI-C247-isliye, HI-C247-agar, HI-C247-ya

\chapter{Silk, Leather, So, If, Or}
\label{ch:hi-silk-leather-so-if-or}
\hlchaptermodality{247}

\begin{chapteropening}
\textbf{By the end of this chapter:} \emph{I can say silk, leather, so, if and or.}

Complete the last lesson of chapter 247: I can say silk, leather, so, if and or.
\end{chapteropening}

\section[\texorpdfstring{\hi{रेशम} (reśam)}{रेशम (reśam)}]{\hi{रेशम} (reśam) — silk}

\label{lesson:HI-C247-resam}



\begin{quote}
Before the new one: say the Hindi for plastic, then the Hindi for cotton wool.
\end{quote}

\subsection*{You'll want to know: \hi{रेशम}}
\textbf{\hi{रेशम}} — \emph{reśam} — \textquotedblleft{}silk\textquotedblright{}.

\begin{cousinweb}
\textbf{\hi{रेशम}} (\emph{reśam}) is a loan, usually traced to Persian \emph{abrēsham}, \textquotedblleft{}silk\textquotedblright{}.
\end{cousinweb}

\begin{grammarlens}[title={What you've built}]
One more word for this chapter.
\end{grammarlens}

\subsection*{Your turn}
\begin{itemize}
\raggedright
  \item \emph{Say it:} \emph{reśam}
  \item \emph{Say it:} \emph{reśam}, once more
  \item \emph{Say it:} say \emph{reśam}
  \item \emph{Recall:} say the Hindi for plastic, then the Hindi for cotton wool, then say \emph{reśam} again
\end{itemize}

\begin{tcolorbox}[breakable,colback=teal!4,colframe=teal!35!black,title={Before you move on}]
What does \hi{रेशम} mean? (\textquotedblleft{}Silk\textquotedblright{}.) Say it once more.
\end{tcolorbox}
\section[\texorpdfstring{\hi{चमड़ा} (cāmṛā)}{चमड़ा (cāmṛā)}]{\hi{चमड़ा} (cāmṛā) — leather}

\label{lesson:HI-C247-camra}



\begin{quote}
Before the new one: say the Hindi for cotton wool, then the Hindi for silk.
\end{quote}

\subsection*{You'll want to know: \hi{चमड़ा}}
\textbf{\hi{चमड़ा}} — \emph{cāmṛā} — \textquotedblleft{}leather\textquotedblright{}.

\begin{cousinweb}
\textbf{\hi{चमड़ा}} (\emph{cāmṛā}) is a tadbhava, a word inherited from Sanskrit and reshaped by sound change. It comes from \emph{carman}, \textquotedblleft{}hide, skin\textquotedblright{} (CDIAL 4701). CDIAL notes an added \emph{-ṛ-}.
\end{cousinweb}

\begin{grammarlens}[title={What you've built}]
One more word for this chapter.
\end{grammarlens}

\subsection*{Your turn}
\begin{itemize}
\raggedright
  \item \emph{Say it:} \emph{cāmṛā}
  \item \emph{Say it:} \emph{cāmṛā}, once more
  \item \emph{Say it:} say \emph{cāmṛā}
  \item \emph{Recall:} say the Hindi for cotton wool, then the Hindi for silk, then say \emph{cāmṛā} again
\end{itemize}

\begin{tcolorbox}[breakable,colback=teal!4,colframe=teal!35!black,title={Before you move on}]
What does \hi{चमड़ा} mean? (\textquotedblleft{}Leather\textquotedblright{}.) Say it once more.
\end{tcolorbox}
\section[\texorpdfstring{\hi{इसलिए} (isliye)}{इसलिए (isliye)}]{\hi{इसलिए} (isliye) — so, therefore}

\label{lesson:HI-C247-isliye}



\begin{quote}
Before the new one: say the Hindi for silk, then the Hindi for leather.
\end{quote}

\subsection*{You'll want to know: \hi{इसलिए}}
\textbf{\hi{इसलिए}} — \emph{isliye} — \textquotedblleft{}so, therefore\textquotedblright{}.

\begin{cousinweb}
\textbf{\hi{इसलिए}} (\emph{isliye}) is \emph{is}, \textquotedblleft{}this\textquotedblright{}, with \emph{liye}, \textquotedblleft{}for\textquotedblright{}: \textquotedblleft{}for this\textquotedblright{}. \emph{liye} is usually explained as a form of \emph{lenā}, \textquotedblleft{}to take\textquotedblright{}, from Sanskrit \emph{labhate} (CDIAL 10948).
\end{cousinweb}

\begin{grammarlens}[title={What you've built}]
One more word for this chapter.
\end{grammarlens}

\subsection*{Your turn}
\begin{itemize}
\raggedright
  \item \emph{Say it:} \emph{isliye}
  \item \emph{Say it:} \emph{isliye}, once more
  \item \emph{Say it:} say \emph{isliye}
  \item \emph{Recall:} say the Hindi for silk, then the Hindi for leather, then say \emph{isliye} again
\end{itemize}

\begin{tcolorbox}[breakable,colback=teal!4,colframe=teal!35!black,title={Before you move on}]
What does \hi{इसलिए} mean? (\textquotedblleft{}So, therefore\textquotedblright{}.) Say it once more.
\end{tcolorbox}
\section[\texorpdfstring{\hi{अगर} (agar)}{अगर (agar)}]{\hi{अगर} (agar) — if}

\label{lesson:HI-C247-agar}



\begin{quote}
Before the new one: say the Hindi for leather, then the Hindi for so.
\end{quote}

\subsection*{You'll want to know: \hi{अगर}}
\textbf{\hi{अगर}} — \emph{agar} — \textquotedblleft{}if\textquotedblright{}.

\begin{cousinweb}
\textbf{\hi{अगर}} (\emph{agar}) is a loan, usually traced to Persian \emph{agar}, \textquotedblleft{}if\textquotedblright{}.
\end{cousinweb}

\begin{grammarlens}[title={What you've built}]
One more word for this chapter.
\end{grammarlens}

\subsection*{Your turn}
\begin{itemize}
\raggedright
  \item \emph{Say it:} \emph{agar}
  \item \emph{Say it:} \emph{agar}, once more
  \item \emph{Say it:} say \emph{agar}
  \item \emph{Recall:} say the Hindi for leather, then the Hindi for so, then say \emph{agar} again
\end{itemize}

\begin{tcolorbox}[breakable,colback=teal!4,colframe=teal!35!black,title={Before you move on}]
What does \hi{अगर} mean? (\textquotedblleft{}If\textquotedblright{}.) Say it once more.
\end{tcolorbox}
\section[\texorpdfstring{\hi{या} (yā)}{या (yā)}]{\hi{या} (yā) — or}

\label{lesson:HI-C247-ya}



\begin{quote}
Before the new one: say the Hindi for so, then the Hindi for if.
\end{quote}

\subsection*{You'll want to know: \hi{या}}
\textbf{\hi{या}} — \emph{yā} — \textquotedblleft{}or\textquotedblright{}.

\begin{cousinweb}
\textbf{\hi{या}} (\emph{yā}) is a loan, usually traced to Persian \emph{yā}, \textquotedblleft{}or\textquotedblright{}.
\end{cousinweb}

\begin{grammarlens}[title={What you've built}]
One more word for this chapter.
\end{grammarlens}

\subsection*{Your turn}
\begin{itemize}
\raggedright
  \item \emph{Say it:} \emph{yā}
  \item \emph{Say it:} \emph{yā}, once more
  \item \emph{Say it:} say \emph{yā}
  \item \emph{Recall:} say the Hindi for so, then the Hindi for if, then say \emph{yā} again
\end{itemize}

\begin{tcolorbox}[breakable,colback=teal!4,colframe=teal!35!black,title={Before you move on}]
What does \hi{या} mean? (\textquotedblleft{}Or\textquotedblright{}.) Say it once more.
\end{tcolorbox}
